\section{Global consequences via finite descent}\label{sec:global-sections}

We now combine the local theorem with established descent and
finite-support results. First we identify the adelic tuples obtained
from global sections. We then
give a criterion that makes these sections geometric and apply it to
abelian targets and modular curves. Finally, we describe further
restrictions on the localization tuples without assuming that criterion.
In this section $K$ is a number field. For a
field $F$ of characteristic zero and a smooth proper geometrically
connected curve $D/F$, write $\operatorname{Sec}(D/F)$ for the
conjugacy classes of continuous splittings of
\[
 1\longrightarrow\pi_1(D_{\bar F})\longrightarrow\pi_1(D)
   \longrightarrow G_F\longrightarrow1,
 \qquad G_F=\operatorname{Gal}(\bar F/F),
\]
where conjugation is by the full geometric \'etale fundamental group.
For a smooth proper geometrically connected curve $C/K$ of genus at
least two, the global section conjecture asserts that the point map
\[
 \kappa_C:C(K)\longrightarrow\operatorname{Sec}(C/K)
\]
is bijective. Thus the section classes considered here are the full
arithmetic \'etale classes, without a birational lifting assumption.

\subsection{Localization and descent}\label{subsec:global-localization}

For a proper $K$-scheme $V$, use the modified adelic set
\[
 V(\mathbb A_K)_\bullet
 =\prod_{v\nmid\infty}V(K_v)
  \times\prod_{v\ \mathrm{real}}\pi_0\bigl(V(K_v)\bigr),
\]
where $\pi_0$ denotes the set of connected components. Complex places
are omitted, as in the corrected convention of Stoll's erratum
\cite[Section~1]{StollFiniteDescentErrata}. For geometrically connected
$C$ the complex component factors in the earlier convention are
singletons, but their omission matters for geometrically disconnected torsors.
A point of $C(\mathbb A_K)_\bullet$ survives a torsor $Y\to C$ under a
finite $K$-group scheme $G$ if it lifts to
$Y^\xi(\mathbb A_K)_\bullet$ for some twist $Y^\xi$, where
$\xi\in H^1(K,G)$.
Write $C(\mathbb A_K)_\bullet^{\mathrm{f\text{-}cov}}$ for the points
surviving every such torsor, and
$C(\mathbb A_K)_\bullet^{\mathrm{f\text{-}ab}}$ for those surviving
every such torsor with $G$ abelian
\cite[Definitions~5.2 and~5.4]{StollFiniteDescent}. We identify $C(K)$ with its
diagonal image. The definitions give
\begin{equation}\label{eq:descent-inclusions}
 C(K)\subseteq C(\mathbb A_K)_\bullet^{\mathrm{f\text{-}cov}}
 \subseteq C(\mathbb A_K)_\bullet^{\mathrm{f\text{-}ab}}.
\end{equation}

The local section theorems define a localization map for every such
curve, without any descent equality hypothesis.
Let $s\in\operatorname{Sec}(C/K)$. At every finite place $v$, its
localization is the section of a unique point $P_v\in C(K_v)$ by
Theorem~\ref{thm:section}. At a real place the section map
\[
 \pi_0\bigl(C(K_v)\bigr)\longrightarrow\operatorname{Sec}(C/K_v)
\]
is bijective by the real section theorem; the genus-at-least-two case
is included in the theorem of Bresciani--Vistoli
\cite[Introduction, Theorem]{BrescianiVistoliReal}. At a complex place
the local Galois group is trivial and the local condition is automatic.
Thus $s$ is a Selmer section: it comes from a local point at every
place. Properness makes these localizations a modified adelic point
$x(s)\in C(\mathbb A_K)_\bullet$. We obtain the canonical map
\[
 \operatorname{loc}:\operatorname{Sec}(C/K)
 \longrightarrow C(\mathbb A_K)_\bullet,
 \qquad s\longmapsto x(s).
\]

\begin{proposition}[The localization image]\label{prop:localization-image}
Let $C/K$ be a smooth proper geometrically connected curve of genus at
least two over a number field. Then
\begin{equation}\label{eq:localization-image}
 \operatorname{loc}\bigl(\operatorname{Sec}(C/K)\bigr)
 =C(\mathbb A_K)_\bullet^{\mathrm{f\text{-}cov}}.
\end{equation}
\end{proposition}

\begin{proof}
Let $s\in\operatorname{Sec}(C/K)$.
Choose a representative point of the determined component at each
real place, and any point at each complex place. Harari--Stix's
Theorem~2.1, with $U=1$ and the omitted set of places $S=\varnothing$,
shows that these localizations survive every torsor under a finite
$K$-group scheme; their Remark~2.2(1) identifies this as finite-cover
descent \cite[preprint version, Theorem~2.1 and Remark~2.2(1)]{HarariStixFiniteDescent}.
Consequently
\[
 x(s)\in C(\mathbb A_K)_\bullet^{\mathrm{f\text{-}cov}}.
\]

Conversely, let $x$ belong to the
finite-cover locus and, by properness, choose an ordinary adelic
representative $P=(P_v)\in C(\mathbb A_K)$. Stoll's cohomological description after
Definition~5.2, together with Definition~5.4, identifies survival with
diagonal global $H^1(K,G)$-membership of the local evaluations for
every torsor under a finite $K$-group scheme $G$
\cite{StollFiniteDescent}. These evaluations factor through the real
components, and complex places impose no lifting condition, in the
corrected convention already used above. Thus every such $P$ satisfies
the ordinary finite-torsor condition. Harari--Stix's Theorem~2.1 with
$U=1$ and $S=\varnothing$, using implication (i)$\Rightarrow$(iii),
gives a global section whose restrictions are the point sections of
$P_v$, up to conjugation by the full geometric fundamental group
\cite[preprint version, Theorem~2.1 and Remark~2.2(1)]{HarariStixFiniteDescent}.
The finite-point and real-component uniqueness used above identify its
localization with $x$, proving~\eqref{eq:localization-image}.
\end{proof}

\subsection{Criteria for point sections}\label{subsec:global-criteria}

Proposition~\ref{prop:localization-image} determines which adelic tuples
come from sections. To recover rational points from the sections
themselves, we also need the finite-support theorem of Stix.

\begin{theorem}[Finite-cover descent criterion]\label{thm:global-descent}
Let $C/K$ be a smooth proper geometrically connected curve of genus at
least two over a number field. If
\begin{equation}\label{eq:finite-cover-equality}
 C(\mathbb A_K)_\bullet^{\mathrm{f\text{-}cov}}=C(K),
\end{equation}
then $\kappa_C:C(K)\to\operatorname{Sec}(C/K)$ is bijective.
\end{theorem}

\begin{proof}
Let $s\in\operatorname{Sec}(C/K)$. By
Proposition~\ref{prop:localization-image} and
\eqref{eq:finite-cover-equality}, there is $a\in C(K)$ such that
$x(s)$ is the diagonal localization of $a$, hence equals the
localization of the point section $s_a$.

At every finite place the point determined by $s$ is therefore $a$.
At each real place choose $a$ itself in the determined component,
and choose $a$ also at each complex place. The resulting ordinary
adelic representative of the localization tuple is supported on
$\{a\}$. Thus the support $Z(s)$ in Stix's terminology, the
Zariski-closed support of this tuple, is finite over $K$.
Since $C$ is proper of genus at least two, it is hyperbolic and has no
cuspidal sections. Stix's finite-support result therefore makes $s$
a point section \cite[arXiv version, Section~3 and Corollary~5]{StixFiniteSupport}.
Finally, if two points give the same global section, their sections
agree after restriction to any one finite completion. Injectivity in
Theorem~\ref{thm:section} then identifies the two points. This proves
injectivity and also identifies the point giving $s$ with $a$.
\end{proof}

The equality in~\eqref{eq:finite-cover-equality} is supplied by several
published criteria. For an abelian variety $A/K$, let
\[
 \operatorname{Sha}(A/K)
 =\ker\left(H^1(K,A)\longrightarrow\prod_v H^1(K_v,A)\right),
\]
and write $\operatorname{Sha}(A/K)_{\mathrm{div}}$ for its maximal
divisible subgroup.

\begin{corollary}[Abelian targets]\label{cor:global-abelian-target}
Let $C/K$ be a smooth proper geometrically connected curve of genus at
least two over a number field. The map $\kappa_C$ is bijective if either
of the following conditions holds:
\begin{enumerate}[label=\textnormal{(\roman*)}]
\item there is a nonconstant $K$-morphism $C\to A$ to an abelian
variety with $A(K)$ finite and $\operatorname{Sha}(A/K)_{\mathrm{div}}=0$;
\item the Jacobian $J$ of $C$ has $J(K)$ finite and
$\operatorname{Sha}(J/K)_{\mathrm{div}}=0$.
\end{enumerate}
\end{corollary}

\begin{proof}
Stoll's Theorem~8.6 gives
$C(\mathbb A_K)_\bullet^{\mathrm{f\text{-}ab}}=C(K)$ under the first
condition, and his Corollary~8.1 gives the same equality under the
second, without requiring a $K$-rational base point on $C$
\cite{StollFiniteDescent}. The corrected covering arguments and the
repair of Proposition~8.5 used by Theorem~8.6 are in
\cite[Sections~2--3]{StollFiniteDescentErrata}; the stated equalities
are retained. The inclusions~\eqref{eq:descent-inclusions} now imply
\eqref{eq:finite-cover-equality}, so Theorem~\ref{thm:global-descent}
applies.
\end{proof}

For an elliptic curve $E/\mathbb Q$ and a prime $p$, put
\[
 s_p(E)=\operatorname{corank}_{\mathbb Z_p}
        \operatorname{Sel}_{p^\infty}(E/\mathbb Q),
\]
where the Selmer group uses the local Kummer conditions at every place.

\begin{corollary}[A Selmer-corank-zero elliptic target]\label{cor:global-selmer-target}
Let $C/\mathbb Q$ be a smooth proper geometrically connected curve of
genus at least two. Suppose that there is a nonconstant
$\mathbb Q$-morphism $C\to E$ to an elliptic curve and that $s_p(E)=0$
for at least one prime $p$. Then $\kappa_C$ is bijective.
\end{corollary}

\begin{proof}
\emergencystretch=2em
The corank-zero case of the companion Selmer converse
\cite[Theorem~1.1]{OpenAISelmerConverse2026} gives
$\operatorname{rank}_{\mathbb Z}E(\mathbb Q)=0$ and finiteness of the
entire group $\operatorname{Sha}(E/\mathbb Q)$. The Mordell--Weil
theorem therefore makes $E(\mathbb Q)$ finite, while finiteness of
$\operatorname{Sha}(E/\mathbb Q)$ makes its maximal divisible subgroup
zero. Corollary~\ref{cor:global-abelian-target}(i) applies.
\end{proof}

\begin{corollary}[Modular curves]\label{cor:global-modular-curves}
For every positive integer $N$ such that the smooth projective
geometrically connected curve $X_0(N)/\mathbb Q$ has genus at least two,
the map $\kappa_{X_0(N)}$ is bijective. The same assertion holds for the
smooth projective geometrically connected curve $X_1(N)/\mathbb Q$
whenever its genus is at least two.
\end{corollary}

\begin{proof}
Stoll's Corollary~8.8 gives
$C(\mathbb A_{\mathbb Q})_\bullet^{\mathrm{f\text{-}ab}}=C(\mathbb Q)$
for each of these projective modular curves of positive genus
\cite{StollFiniteDescent}, with the corrected conventions of
\cite{StollFiniteDescentErrata}. Apply~\eqref{eq:descent-inclusions}
and Theorem~\ref{thm:global-descent}.
\end{proof}

The general implication retains the finite-cover equality as a genuine
arithmetic hypothesis; local point-theoreticity alone does not supply
it. The elliptic-target specialization uses only Selmer corank zero over
$\mathbb Q$, while the modular-curve cases come directly from Stoll's
descent result and do not depend on that specialization.

\subsection{Further restrictions on localization}
\label{subsec:global-restrictions}

We return to an arbitrary smooth proper geometrically connected curve
$C/K$ of genus at least two over a number field, without
assuming~\eqref{eq:finite-cover-equality}. The localization image in
Proposition~\ref{prop:localization-image} satisfies additional descent
conditions and, under a restriction on $K$, has finite image at each
finite place.

Let $\rho:C(\mathbb A_K)\to C(\mathbb A_K)_\bullet$ be the natural
projection. Write $D_{\mathrm{lin}}(C)$ for the ordinary adelic descent
set: $P$ belongs to it if, for every $C$-torsor $T\to C$ under a linear
algebraic $K$-group $G$, it lifts to $T^\xi(\mathbb A_K)$ for some
global twist $\xi\in H^1(K,G)$. Harari--Stix's Theorem~4.1 sends the
finite-torsor condition established in
Proposition~\ref{prop:localization-image} to this linear descent
condition, while their Corollary~3.1, applied to the geometric
abelianization of the interpolating section, gives Brauer--Manin
membership
\cite[preprint version, Corollary~3.1 and Theorem~4.1]{HarariStixFiniteDescent}.
Consequently
\begin{equation}\label{eq:localization-obstructions}
 \rho^{-1}\!\left(\operatorname{loc}\bigl(\operatorname{Sec}(C/K)\bigr)\right)
 \subseteq D_{\mathrm{lin}}(C)\cap C(\mathbb A_K)^{\operatorname{Br}}.
\end{equation}
This applies to every ordinary adelic representative of a modified
localization tuple; the lifting condition uses ordinary adeles also
when $T$ is nonproper.

Finally, if $K$ contains no CM subfield, Betts--Stix's Theorem~A,
applied to the sections shown to be Selmer in
Section~\ref{subsec:global-localization}, gives
\[
 \#\operatorname{loc}_v\bigl(\operatorname{Sec}(C/K)\bigr)<\infty,
 \qquad
 \operatorname{loc}_v=\operatorname{pr}_v\circ\operatorname{loc},
\]
for every finite place $v$, with image in $C(K_v)$
\cite[arXiv version, Definition~1.1 and Theorem~A]{BettsStix}.
This asserts neither injectivity of localization nor finiteness of the
global section set, and supplies no effective or uniform bound.
